L'obiettivo inferenziale è stabilire quali proprietà delle due memorie frazionarie possano essere ricostruite dalla sola concentrazione del tracciante. Nel sistema~\eqref{eq:coupled-system}, $\alpha$ e $D$ agiscono direttamente sull'equazione di trasporto, mentre $\beta$ ed $E$ modificano il campo di pressione e, attraverso la legge di Darcy, il trasporto avvettivo. Tale struttura rende distinti tre problemi: dipendenza delle osservazioni dai parametri, identificabilità locale della mappa e stabilità della loro ricostruzione in presenza di rumore o discrepanza numerica. La formulazione che segue tiene separati questi livelli.

\subsection{Spazio parametrico e operatore osservazionale}

Consideriamo innanzitutto quattro parametri adimensionali e il riferimento nominale
\begin{equation}
 \tht=(\alpha,\beta,D,E)^{\mathsf T},\qquad
 \tht^\dagger=(0.70,0.45,0.08,2.50)^{\mathsf T},
 \label{eq:theta-physical}
\end{equation}
con insieme ammissibile
\begin{equation}
 \Theta_{\rm inv}=[0.15,0.90]^2\times[0.01,0.50]\times[0.50,6].
 \label{eq:theta-inverse-box}
\end{equation}
Negli esperimenti anatomici, $D$ ed $E$ sono coefficienti nominali moltiplicati per i campi regionali prescritti nella Sezione~\ref{sec:forward}; tali campi restano noti nell'inversione. Per sondare la compensazione con la risposta idraulica e con la memoria viscosa, il vettore esteso è $\boldsymbol\vartheta=(\tht^{\mathsf T},K_h,A_v)^{\mathsf T}$, con $K_h\in[0.10,0.40]$ e $A_v\in[0.225,0.90]$. In questo protocollo $\tau_v=1$, sicché $A_v=\eta_v$ può variare indipendentemente da $\beta$. La dimensione dell'incognita è dunque $k\in\{4,6\}$ e non comprende una ricostruzione libera dei moltiplicatori regionali.

Siano $\mathcal S=\{\mathbf x_i\}_{i=1}^{n_s}$ i sensori, $\mathcal T=\{t_j\}_{j=1}^{n_t}$ gli istanti osservati e $N_{\rm obs}=n_s n_t$. Distinguiamo con $\mathsf d=(\mathcal M_h,\mathcal T_h)$ la discretizzazione spaziale e temporale, per non attribuire al solo passo spaziale tutte le differenze fra due operatori. La mappa discreta di osservazione è
\begin{equation}
 \begin{aligned}
 \Gmap_{\mathsf d,\mathcal S}(\boldsymbol\vartheta)
  &=\mathcal H_{\mathsf d,\mathcal S}\,c_{\mathsf d}(\boldsymbol\vartheta)
    \in\R^{N_{\rm obs}},\\
 \mathbf y^\dagger
  &=\Gmap_{\mathsf d_{\rm tar},\mathcal S}(\boldsymbol\vartheta^\dagger),
  \qquad \mathbf y=\mathbf y^\dagger+\boldsymbol\epsilon .
 \end{aligned}
 \label{eq:observation}
\end{equation}
L'operatore $\mathcal H$ estrae la concentrazione in punti fisici e tempi prestabiliti; l'interpolazione temporale fra due nodi consecutivi è lineare. Sul quadrato si usano $36$ sensori su una griglia $6\times6$ e cinque tempi ($N_{\rm obs}=180$), con campionamento spaziale \texttt{griddata} sulla triangolazione dei nodi del tracciante. Nella prima diagnosi anatomica sono impiegati $36$ sensori focali e l'interpolazione locale $Q_1$ della cella quadrilatera. Nei confronti anatomici regionali si usano invece due disegni \emph{prefissati} da $16$ sensori ($N_{\rm obs}=80$): uno bilanciato fra le quattro regioni e uno concentrato nell'edema. I tempi sono $t_j=0.12j$ per $j=1,\ldots,5$. La scelta dei sensori non è il risultato di un'ottimizzazione sul Jacobiano.

Per rendere confrontabili le ampiezze parametrizzate, si introducono coordinate $\mathbf z\in[0,1]^k$. Per una componente $u$ con estremi $u_-,u_+$ si pone
\begin{equation}
\begin{aligned}
 z_u&=\frac{g_u(u)-g_u(u_-)}{g_u(u_+)-g_u(u_-)},\\[-1pt]
 g_u(u)&=\begin{cases}
   u, & u\in\{\alpha,\beta\},\\
   \log u, & u\in\{D,E,K_h,A_v\}.
 \end{cases}
\end{aligned}
 \label{eq:normalized-parameters}
\end{equation}
La scelta logaritmica preserva la positività dei coefficienti e rende una perturbazione unitaria in $z_u$ equivalente a un rapporto moltiplicativo prefissato.

\subsection{Problema ai minimi quadrati, rumore e regolarizzazione}

Per una griglia inversa $\mathsf d_{\rm inv}$, definiamo $\mathbf r(\mathbf z)=\Gmap_{\mathsf d_{\rm inv},\mathcal S}(\boldsymbol\vartheta(\mathbf z))-\mathbf y$. La ricostruzione risolve il problema vincolato
\begin{equation}
 \Jfun(\mathbf z)=
 \underbrace{\tfrac12\|W\mathbf r(\mathbf z)\|_2^2}_{J_{\rm data}}
 +\underbrace{\tfrac\lambda2\|\mathbf z-\mathbf z_0\|_2^2}_{J_{\rm reg}},
 \qquad \lambda\ge0 .
 \label{eq:objective}
\end{equation}
\begin{equation}
 \widehat{\mathbf z}\in\operatorname*{arg\,min}_{\mathbf z\in[0,1]^k}
 \Jfun(\mathbf z),\qquad k\in\{4,6\}.
 \label{eq:inverse-min}
\end{equation}
Se il rumore è gaussiano con covarianza positiva definita $\Sigma$, la scelta statistica $W^{\mathsf T}W=\Sigma^{-1}$ produce, a costanti additive irrilevanti, il negativo della log-verosimiglianza. Altre normalizzazioni di $W$ definiscono invece convenzioni deterministiche o modificano il peso relativo della penalizzazione.

\paragraph{Quadrato sintetico.}
Con $s_{\rm ref}=\|\mathbf y^\dagger\|_\infty$, la deviazione standard del rumore è $\sigma=\eta_\infty s_{\rm ref}$. Si distinguono la normalizzazione deterministica A e due normalizzazioni dei residui rumorosi B:
\begin{equation}
\begin{aligned}
 W_A&=I/(\sqrt{N_{\rm obs}}\,s_{\rm ref}),\\
 W_{B,\rm mean}&=I/(\sqrt{N_{\rm obs}}\,\sigma),\qquad
 W_{B,\rm cov}=I/\sigma .
\end{aligned}
\label{eq:weight-matrices}
\end{equation}
Per lo stesso residuo, con $\sigma=\eta_\infty s_{\rm ref}$, vale esattamente
\begin{equation}
 J_{\rm data}^{B,\rm mean}=\eta_\infty^{-2}J_{\rm data}^A,
 \qquad J_{\rm data}^{B,\rm cov}=N_{\rm obs}\eta_\infty^{-2}J_{\rm data}^A.
 \label{eq:weight-ratios}
\end{equation}
Ne consegue che conservare lo stesso valore numerico di $\lambda$ mentre si modifica $W$ non conserva il rapporto fra dati e prior. I twin puliti utilizzano $W_A$, $\lambda_A=10^{-4}$ e $\mathbf z_0=(0.5,0.5,0.5,0.5)^{\mathsf T}$; i controlli rumorosi e la conversione della penalizzazione sono precisati nella Sezione~\ref{sec:objective-protocol}.

\paragraph{Geometria anatomica.}
Per inversioni e profili anatomici si pone $\lambda=0$. La scala assoluta comune ai due disegni è
\[
 \sigma_0=0.05\,\operatorname{RMS}(\mathbf y^\dagger_{3\,\mathrm{mm},\rm bil}),
 \qquad \operatorname{RMS}(\mathbf y)=\|\mathbf y\|_2/\sqrt{N_{\rm obs}}.
\]
I profili sui dati puliti usano $W=I/\sigma_0$ come riferimento di informazione. Nelle inversioni rumorose si genera $\boldsymbol\epsilon\sim\mathcal N(0,\sigma_\eta^2I)$ con
$\sigma_\eta=\eta_{\rm RMS}\operatorname{RMS}(\mathbf y^\dagger_{3\,\mathrm{mm},\rm bil})$ e si pesa il residuo mediante $W=I/\sigma_\eta$. Questa scala resta la stessa per sensori bilanciati ed edema a uno stesso livello nominale.

È utile distinguere il parametro nominale dalla quantità realizzata su un campione,
$\eta_{\rm eff}=\|\boldsymbol\epsilon\|_2/\|\mathbf y^\dagger\|_2$.
Sul quadrato il rapporto RMS \emph{nominale} vale $\eta_\infty\|\mathbf y^\dagger\|_\infty/\operatorname{RMS}(\mathbf y^\dagger)$; sul disegno anatomico bilanciato coincide con $\eta_{\rm RMS}$. Nell'edema è inoltre moltiplicato per $\operatorname{RMS}(\mathbf y^\dagger_{\rm bil})/\operatorname{RMS}(\mathbf y^\dagger_{\rm edema})$. Le percentuali nominali dei diversi domini non sono quindi direttamente confrontabili. Il numero di condizionamento di $J$ resta invariato sotto una riscalatura scalare uniforme delle righe, ma in generale non sotto cambiamenti dell'operatore osservazionale o della pesatura relativa.

Se $\mathsf d_{\rm tar}\ne\mathsf d_{\rm inv}$, anche senza rumore si ha
\begin{equation}
 \mathbf r(\mathbf z^\dagger)
 =\Gmap_{\mathsf d_{\rm inv},\mathcal S}(\boldsymbol\vartheta^\dagger)
  -\Gmap_{\mathsf d_{\rm tar},\mathcal S}(\boldsymbol\vartheta^\dagger),
 \label{eq:inverse-model-discrepancy}
\end{equation}
una discrepanza deterministica distinta dall'errore osservativo. Essa può essere dovuta all'accuratezza spazio-temporale, alla geometria approssimata o alla rappresentazione materiale: non coincide necessariamente con un errore spaziale FEM isolato~\citep{kaipio2007discretization}.

\subsection{Canale meccanica--flusso--tracciante e regime di eccitazione}

\refstepcounter{theorem}\label{prop:pv-zero-flow}
\paragraph{Proposizione~\thetheorem\ (invisibilità meccanica a flusso nullo).}
Si fissino $\alpha,D,\phi$, le condizioni iniziali, il termine sorgente e i dati al bordo del tracciante. Se per ogni coppia $(E,\beta)$ in un intorno del riferimento il flusso $\mathbf q$ è identicamente nullo e il problema del trasporto ha soluzione unica, allora la concentrazione è indipendente da $(E,\beta)$ in tale intorno. Infatti l'equazione del trasporto si riduce a
\begin{equation}
 \phi\CD{\alpha}c-\nabla\!\cdot(\phi D\nabla c)=s_c ,
 \label{eq:zero-flow-tracer}
\end{equation}
che non contiene $E$ né $\beta$. Ne segue che, ove esistano, le corrispondenti colonne del Jacobiano osservativo sono nulle. Lo stesso argomento vale per la mappa discreta se il suo flusso è nullo in quell'intorno parametrico.

Per una perturbazione dei parametri meccanici, con $K_h$ fissato, risulta $\delta\mathbf q=-\mathbb K_h\nabla\delta p$. La linearizzazione del trasporto attorno a una traiettoria $(c,\mathbf q)$ dà
\begin{equation}
\begin{aligned}
 \phi\CD{\alpha}\delta c+\nabla\!\cdot(\delta c\,\mathbf q)
  -\nabla\!\cdot(\phi D\nabla\delta c)
  &=-\nabla\!\cdot(c\,\delta\mathbf q).
\end{aligned}
\label{eq:mech-tracer-variation}
\end{equation}
La sorgente destra descrive il canale osservabile mediante il quale la memoria del solido può influire sulla concentrazione. Per una storia $u$ fissata sufficientemente regolare, la sensibilità \emph{esplicita} dell'operatore di Caputo rispetto all'ordine è
\[
\begin{aligned}
 &\left.\partial_\gamma\bigl(\CD{\gamma}u\bigr)(t)\right|_{u\,\mathrm{fisso}}\\
 &\quad=\int_0^t k_\gamma(t-s)u'(s)
 [\psi(1-\gamma)-\log(t-s)]\,\dd s .
\end{aligned}
\]
Le sensibilità complete includono anche la dipendenza implicita delle soluzioni da $\gamma$. Nella memoria meccanica, mantenendo fissi $\eta_v$ e $\tau_v$, interviene inoltre $\partial_\beta A_v=A_v\log\tau_v$, come previsto dalla legge costitutiva della Sezione~\ref{sec:forward}.

\paragraph{Scala delle direzioni deboli.}
Poniamo $\varepsilon=\kappa_q\ge0$ e ordiniamo le colonne del Jacobiano pesato come $J(\varepsilon)=[A(\varepsilon),B(\varepsilon)]$, dove $A$ corrisponde a $(\alpha,D)$ e $B$ a $(\beta,E)$. Assumiamo che la mappa \emph{discreta}, in un intorno di $\varepsilon=0$, sia sufficientemente regolare e soddisfi
$A(\varepsilon)=A_0+O(\varepsilon)$, $B(\varepsilon)=\varepsilon B_1+O(\varepsilon^2)$, con $A_0$ a rango colonnare due. Definendo $P_0=A_0A_0^+$ e $B_\perp=(I-P_0)B_1$, la perturbazione dei valori singolari conduce a
\begin{equation}
 \sigma_{2+j}(J(\varepsilon))
 =\varepsilon\,\sigma_j(B_\perp)+o(\varepsilon),\qquad j=1,2,
 \label{eq:weak-flow-singular}
\end{equation}
con i due valori singolari di $B_\perp$ ordinati in senso decrescente. Per giustificare la relazione, si proiettano le ultime due colonne sul complemento dello spazio generato da $A_0$: la loro componente interna a tale spazio si elimina al primo ordine mediante trasformazioni di colonna prossime all'identità, mentre il blocco ortogonale residuo è $\varepsilon B_\perp+O(\varepsilon^2)$. Se $B_\perp$ ha rango due, $J(\varepsilon)$ ha rango quattro per $\varepsilon>0$ sufficientemente piccolo. Se perde rango, almeno una direzione può emergere a ordine superiore.

\subsection{Identificabilità locale e compensazione parametrica}

Fissati la griglia, i sensori, il riferimento e la scala del residuo, definiamo
\[
 J_{\rm raw}=\partial_{\mathbf z}\Gmap_{\mathsf d,\mathcal S}(\boldsymbol\vartheta(\mathbf z^\dagger)),
 \qquad J=WJ_{\rm raw},\qquad F=J^{\mathsf T}J.
\]
Quando $W=\Sigma^{-1/2}$, $F$ è la matrice di informazione di Fisher del modello gaussiano linearizzato. Con $W=I$, $F$ è invece una matrice di Gram delle sensitività e non va interpretata come informazione statistica calibrata. Lo spettro singolare, il rango e $\kappa_2(J)=\sigma_{\max}(J)/\sigma_{\min}(J)$ descrivono proprietà locali della mappa in coordinate normalizzate.

\refstepcounter{theorem}\label{prop:pv-iniettivita}
\paragraph{Proposizione~\thetheorem\ (iniettività locale della mappa discreta).}
Sia $\Gmap_{\mathsf d,\mathcal S}$ di classe $C^1$ nell'intorno di un punto \emph{interno} dello spazio parametrico. Se il Jacobiano ha rango colonnare pieno, la mappa è localmente iniettiva. Infatti esiste una selezione di $k$ componenti osservate il cui Jacobiano quadrato è invertibile; il teorema della funzione inversa applicato a tale selezione dà un'inversa locale. Per $W$ invertibile il rango di $WJ_{\rm raw}$ coincide con quello del Jacobiano grezzo. Questa conclusione è locale e non esclude soluzioni inverse multiple lontane dal riferimento.

In regime lineare, la pseudoinversa di $J$ amplifica al massimo un errore osservativo pesato di norma unitaria di un fattore $1/\sigma_{\min}(J)$, se il rango è pieno. Per il modello a sei parametri scriviamo $J=[J_\theta\;J_m]$, dove $J_m$ raccoglie le colonne di $(K_h,A_v)$, e definiamo
\[
 \mathcal I_{\theta\mid m}=J_\theta^{\mathsf T}(I-J_mJ_m^+)J_\theta .
\]
Questo complemento geometrico rimuove la parte delle sensibilità principali che gli altri due coefficienti possono riprodurre al primo ordine. In particolare, per $j_\beta$ colonna di $J$ relativa a $z_\beta$, si pone
\begin{equation}
 r_\beta=(I-P_{-\beta})j_\beta,\qquad
 s_{\beta\mid-\beta}=\|r_\beta\|_2,\qquad
 \mathcal I_{\beta\mid-\beta}=s_{\beta\mid-\beta}^2 ,
 \label{eq:conditional-beta-info}
\end{equation}
dove $P_{-\beta}$ è la proiezione ortogonale sulle colonne complementari. Tale quantità è strettamente connessa alla varianza del singolo parametro e alla progettazione sperimentale $c$-ottimale~\citep{pukelsheim2006optimal}. Se $W=\Sigma^{-1/2}$, il rango è pieno, non vi sono vincoli attivi e il modello linearizzato è appropriato, allora
\[
 \operatorname{Cov}_{\rm loc}(\widehat{\mathbf z})=(J^{\mathsf T}J)^{-1},
 \qquad \operatorname{sd}_{\rm loc}(\widehat\beta)
 =\frac{\beta_+-\beta_-}{s_{\beta\mid-\beta}}.
\]
Quando $r_\beta=0$, la variazione infinitesima di $\beta$ è indistinguibile da una compensazione delle altre colonne. La covarianza locale non descrive, senza ulteriori controlli, la distribuzione di uno stimatore non lineare vincolato, soprattutto in prossimità dei bordi.

Per investigare compensazioni non lineari fissiamo $\beta=b_*$ e minimizziamo il solo termine dati sugli altri coefficienti:
\[
 P_\beta(b_*)=\inf_{\substack{\mathbf z\in[0,1]^k\\\beta(\mathbf z)=b_*}} J_{\rm data}(\mathbf z).
\]
I profili \emph{numerici} ottenuti con TRF forniscono minimi locali candidati, ossia limiti superiori al profilo globalmente minimizzato. Sul quadrato si calcola anche l'analogo profilo del costo regolarizzato. In anatomia, con dati puliti generati e invertiti mediante la medesima mappa, il costo al riferimento è nullo. Il contrasto $\Delta\chi^2(b_*)=2[P_\beta(b_*)-P_\beta(\beta^\dagger)]$ ammette la previsione locale
\[
 \Delta\chi^2_{\rm lin}(b_*)
 =s_{\beta\mid-\beta}^2
 \left(\frac{b_*-\beta^\dagger}{\beta_+-\beta_-}\right)^2 .
\]
Il confronto fra i punti riottimizzati e tale parabola misura la compensazione rispetto alla linearizzazione, non garantisce intervalli di confidenza da tre soli nodi~\citep{raue2009structural,brun2001practical}.

\subsection{Discrepanza numerica e previsione locale del bias}

Per isolare l'effetto di discretizzazioni differenti occorrono \emph{la stessa mappa costitutiva}, i medesimi sensori fisici, gli stessi tempi osservati e la stessa pesatura. Siano $\Gmap_c$ la mappa inversa grossolana e $\Gmap_f$ quella del target fine, entrambe valutate al riferimento $\mathbf z^\dagger$. Definiamo
\[
 e_h=W[\Gmap_f(\mathbf z^\dagger)-\Gmap_c(\mathbf z^\dagger)],
 \qquad J_c=W\,\partial_{\mathbf z}\Gmap_c(\mathbf z^\dagger).
\]
Per il problema lineare \emph{non regolarizzato} e con soluzione interna, la perturbazione prevista risolve $\min_{\delta\mathbf z}\|J_c\delta\mathbf z-e_h\|_2^2$. Separando $z_\beta$ dagli altri parametri e proiettando fuori le loro sensibilità, la formula di eliminazione ai minimi quadrati fornisce
\begin{equation}
\begin{aligned}
 \widehat{\delta z_\beta}^{\,\rm lin}
 &=\frac{\langle r_\beta,e_h\rangle}{\|r_\beta\|_2^2},\\
 \widehat{\delta\beta}^{\,\rm lin}
 &=(\beta_+-\beta_-)\widehat{\delta z_\beta}^{\,\rm lin} .
\end{aligned}
\label{eq:beta-discretization-bias}
\end{equation}
È necessario $r_\beta\ne0$ e una pseudoinversa sufficientemente stabile per le altre colonne. Se la previsione supera i vincoli, è appropriato esaminare anche il problema linearizzato vincolato; se il residuo numerico è grande, la non linearità può invalidare la previsione. Il confronto fra le sole norme $\|e_h\|_2$ e $\|\mathbf y(\beta_1)-\mathbf y(\beta_2)\|_2$ non ne determina né il segno né l'ampiezza.

I valori di riferimento della distribuzione $\chi_1^2$, come il quantile $3.841$, sono interpretabili come soglie asintotiche soltanto sotto le ipotesi statistiche appropriate.
